%! Zeros of Polynomials: the Rational Root Theorem — Unit 3, Lesson 3.4 — Group Activity (Answer Key)
\documentclass[10pt]{article}
\usepackage{algebra2-article}
\usepackage{algebra2-key}   % pulls in -boxes; do NOT also load -boxes
\newcommand{\TallMath}[1]{$\displaystyle #1\rule[-1.4em]{0pt}{3.2em}$}

\newcommand{\lgrid}[4]{%
  \draw[step=1, linegray!40, very thin] (#1,#3) grid (#2,#4);
  \draw[->, semithick, charcoal] (#1,0)--(#2,0) node[below right, font=\scriptsize]{$x$};
  \draw[->, semithick, charcoal] (0,#3)--(0,#4) node[left, font=\scriptsize]{$y$};}

\begin{document}

\pageheader{Unit 3, Lesson 3.4}{Group Activity --- Answer Key}
\namepartnerperiod

\vspace{0.08in}

\begin{headlinebox}{forestbg}
\small\textbf{Suspect List.} Every problem today runs on the same four moves: \textbf{list} the
candidates $\pm\frac{p}{q}$, \textbf{test} them, \textbf{divide} out the winner, and \textbf{finish}
with the depressed polynomial. Two habits keep you honest: the constant term goes on \emph{top} and
the leading coefficient on the \emph{bottom}, and every candidate comes with a $\pm$. Remember what
the list is --- a list of \emph{suspects}, not convictions. Some of them will not be zeros at all.
\end{headlinebox}

\vspace{0.06in}

\begin{tcolorbox}[colback=white, colframe=black!40, title=\textbf{Tier R --- Remediate},
    fonttitle=\bfseries, arc=1.5mm, left=3mm, right=3mm, top=2mm, bottom=2mm, breakable]
\begin{enumerate}[leftmargin=1.7em, itemsep=9pt]
  \item \textbf{Build one list, step by step.} $f(x) = x^3 - 2x^2 - 5x + 6$

    \vspace{2pt}
    The constant term is \ans{$6$}, so the numerators $p$ are \ans{$\pm 1, \pm 2, \pm 3, \pm 6$}

    \vspace{2pt}
    The leading coefficient is \ans{$1$}, so the denominators $q$ are \ans{$\pm 1$}

    \vspace{2pt}
    Candidate list $\dfrac{p}{q}$: \ans{$\pm 1,\ \pm 2,\ \pm 3,\ \pm 6$}

    \vspace{3pt}
    Now test two of them: \; $f(1) = $ \ans{$0$} \qquad $f(2) = $ \ans{$-4$}

    \vspace{2pt}
    Which one is a zero? \ans{$1$} \; So \ans{$(x-1)$} is a factor of $f$.

    \vspace{3pt}
    Divide it out:
    \begin{center}
    \renewcommand{\arraystretch}{1.6}
    \begin{tabular}{r|C{1.4cm}C{1.4cm}C{1.4cm}C{1.4cm}}
    \ans{$1$} & $1$ & $-2$ & $-5$ & $6$ \\
        &     & \ans{$1$} & \ans{$-1$} & \ans{$-6$} \\ \cline{2-5}
        & \ans{$1$} & \ans{$-1$} & \ans{$-6$} & \ans{$0$} \\
    \end{tabular}
    \end{center}

    \vspace{2pt}
    Depressed polynomial: \ans{$x^2 - x - 6$} \; Factor it: \ans{$(x-3)(x+2)$}

    \vspace{2pt}
    $f(x) = $ \ans{$(x-1)(x-3)(x+2)$} \qquad The three zeros are \ans{$1,\ 3,\ -2$}

  \item \textbf{Same four moves, on your own.} $g(x) = x^3 + 2x^2 - x - 2$

    \vspace{2pt}
    Candidates: \ans{$\pm 1,\ \pm 2$} \qquad A zero that works: $r = $ \ans{$1$}

    \vspace{2pt}
    Depressed polynomial: \ans{$x^2 + 3x + 2$} \qquad $g(x) = $ \ans{$(x-1)(x+1)(x+2)$}

    \vspace{2pt}
    Zeros: \ans{$1,\ -1,\ -2$}

    \vspace{2pt}
    In Lesson 3.2 you could have factored this one by \emph{grouping}. Show that route in one line:
    \ans{$x^2(x+2) - 1(x+2) = (x+2)(x^2-1) = (x+2)(x-1)(x+1)$}

  \item \textbf{When the leading coefficient is not 1.} $h(x) = 2x^3 + x^2 - 5x + 2$

    \vspace{2pt}
    Numerators $p$ (divisors of \ans{$2$}): \ans{$\pm 1, \pm 2$} \qquad
    Denominators $q$ (divisors of \ans{$2$}): \ans{$\pm 1, \pm 2$}

    \vspace{2pt}
    Candidate list: \ans{$\pm 1,\ \pm 2,\ \pm\tfrac12$} \; How many candidates are there? \ans{$6$}

    \vspace{2pt}
    Test $r = 1$: $h(1) = $ \ans{$0$} \; Is $(x-1)$ a factor? \ans{Yes}
\end{enumerate}
\end{tcolorbox}

\begin{tcolorbox}[colback=white, colframe=black!40, title=\textbf{Tier A --- Approaching Proficiency},
    fonttitle=\bfseries, arc=1.5mm, left=3mm, right=3mm, top=2mm, bottom=2mm, breakable]
\begin{enumerate}[leftmargin=1.7em, itemsep=10pt]
  \item \textbf{A fractional zero.} Solve $2x^3 - x^2 - 8x + 4 = 0$.

    \vspace{2pt}
    Candidates: \ans{$\pm 1,\ \pm 2,\ \pm 4,\ \pm\tfrac12$}

    \vspace{2pt}
    A zero that works: $r = $ \ans{$2$} \quad Depressed polynomial: \ans{$2x^2 + 3x - 2$}

    \vspace{2pt}
    Completely factored: \ans{$(x-2)(2x-1)(x+2)$} \quad Solutions: \ans{$x = -2,\ \tfrac12,\ 2$}

    \vspace{2pt}
    This one also factors by grouping. Do it as a check:
    \ans{$x^2(2x-1) - 4(2x-1) = (2x-1)(x^2-4) = (2x-1)(x-2)(x+2)$}

  \item \textbf{Let the graph do the work.} Below is $f(x) = 2x^3 + 3x^2 - 11x - 6$.
    (Each vertical gridline is $4$ units.)
    \begin{center}
    \begin{tikzpicture}[scale=0.5]
      \lgrid{-4}{3}{-4}{4}
      \begin{scope}
        \clip (-4,-4) rectangle (3,4);
        \draw[very thick, forest, domain=-3.4:2.4, samples=150, smooth]
          plot (\x,{(2*(\x)*(\x)*(\x) + 3*(\x)*(\x) - 11*(\x) - 6)/4});
      \end{scope}
      \foreach \x in {-3,-2,-1,1,2}
        \draw[charcoal] (\x,-0.15)--(\x,0.15) node[below=1pt, font=\tiny]{$\x$};
    \end{tikzpicture}
    \end{center}

    \vspace{2pt}
    (a) Write the full candidate list: \ans{$\pm 1,\ \pm 2,\ \pm 3,\ \pm 6,\ \pm\tfrac12,\ \pm\tfrac32$}

    \vspace{2pt}
    (b) The graph crosses the $x$-axis three times. Two of the crossings are at integers ---
    read them: \ans{$-3$ and $2$}

    \vspace{2pt}
    (c) The third crossing is between $x = -1$ and $x = 0$. Which candidates from your list live in
    that interval? \ans{only $-\tfrac12$}

    \vspace{2pt}
    (d) Test the one you think it is, then write $f$ completely factored:
    \ans{$f(-\tfrac12)=0$; \; $f(x) = (x+3)(2x+1)(x-2)$}

    \vspace{2pt}
    (e) All three solutions: \ans{$x = -3,\ -\tfrac12,\ 2$} \; How many synthetic divisions did the graph save you?
    \ans{The list holds $12$ candidates; the graph pointed straight at the three that matter, so
    nothing had to be tested blindly.}

  \item \textbf{Degree 4 --- two trips through the workflow.} Solve $x^4 - x^3 - 7x^2 + x + 6 = 0$.

    \vspace{2pt}
    Candidates: \ans{$\pm 1,\ \pm 2,\ \pm 3,\ \pm 6$}

    \vspace{2pt}
    First zero $r_1 = $ \ans{$1$} \; Depressed polynomial: \ans{$x^3 - 7x - 6$}

    \vspace{2pt}
    Now start over on \emph{that} polynomial. Second zero $r_2 = $ \ans{$-1$} \;
    Next depressed polynomial: \ans{$x^2 - x - 6$}

    \vspace{2pt}
    Factor what is left: \ans{$(x-3)(x+2)$} \quad Completely factored: \ans{$(x-1)(x+1)(x-3)(x+2)$}

    \vspace{2pt}
    Solutions: \ans{$x = -2,\ -1,\ 1,\ 3$} \; A degree-$4$ polynomial has at most \ans{$4$} zeros --- did you
    find them all? \ans{Yes}

  \item \textbf{Error analysis.} Two students, two different misunderstandings.

    \vspace{2pt}
    (a) For $3x^3 - 2x^2 + x - 4$, a student writes the candidate list $1, 2, 4, \tfrac13, \tfrac23,
    \tfrac43$. What is missing, and how many candidates should there be?
    \ans{The negatives --- every candidate needs a $\pm$. There should be $12$, not $6$.}

    \vspace{2pt}
    (b) A student says: ``The candidates for $k(x) = x^3 - 3x - 1$ are $\pm 1$, so one of them has
    to be a zero.'' Test both: $k(1) = $ \ans{$-3$} \quad $k(-1) = $ \ans{$1$}

    \vspace{2pt}
    What is wrong with the student's reasoning? \ans{The theorem lists \emph{possible} rational
    zeros. It promises that any rational zero is on the list --- not that the list contains one.
    This $k$ has no rational zeros at all.}
\end{enumerate}
\end{tcolorbox}

\begin{tcolorbox}[colback=white, colframe=black!40, title=\textbf{Tier E --- Extension},
    fonttitle=\bfseries, arc=1.5mm, left=3mm, right=3mm, top=2mm, bottom=2mm, breakable]
\textbf{Part 1 --- Prove the easy half.} Show that if an \emph{integer} $r$ is a zero of
$f(x) = x^3 + bx^2 + cx + d$ (with $b$, $c$, $d$ integers), then $r$ must divide $d$.
\begin{enumerate}[label=(\alph*), leftmargin=1.8em, itemsep=6pt]
  \item Because $r$ is a zero, $f(r) = 0$, so \; $r^3 + br^2 + cr + d = $ \ans{$0$}
  \item Move $d$ to the other side and factor $r$ out of the three remaining terms:
        $d = -r\big($ \ans{$r^2 + br + c$} $\big)$
  \item The quantity in parentheses is an \ans{integer}, so $d$ is $r$ times an integer, which is
        exactly what ``\,$r$ divides $d$\,'' means. \; $\blacksquare$
  \item In one sentence, why does this argument \emph{not} work when the leading coefficient is
        something other than $1$? \ansline{With a leading coefficient $a_n \ne 1$ the zero need not be an integer at all --- it can be a fraction $\frac{p}{q}$, and clearing that denominator is what brings the leading coefficient into the statement.}
\end{enumerate}

\vspace{5pt}
\textbf{Part 2 --- Run the theorem backward.} Build a polynomial to order.
\begin{enumerate}[label=(\alph*), leftmargin=1.8em, itemsep=6pt]
  \item You want a cubic with integer coefficients whose candidate list is exactly
        $\pm 1, \pm\tfrac12, \pm\tfrac13, \pm\tfrac16$. What must its constant term be?
        \ans{$\pm 1$} \; Its leading coefficient? \ans{$\pm 6$}
  \item Now build one with zeros $\tfrac12$, $\tfrac13$, and $-1$. Start from the factors
        $(2x-1)$, $(3x-1)$, $($\ans{$x+1$}$)$ and multiply them out:
        $f(x) = $ \ans{$6x^3 + x^2 - 4x + 1$}
  \item Check your answer against (a): does its constant term and leading coefficient give the list
        you wanted? \ans{Yes --- constant $1$, leading coefficient $6$}
\end{enumerate}

\vspace{5pt}
\textbf{Part 3 --- A polynomial the theorem cannot crack.} Let $k(x) = x^3 - 3x - 1$.
\begin{enumerate}[label=(\alph*), leftmargin=1.8em, itemsep=6pt]
  \item Candidate list: \ans{$\pm 1$} \; Test every candidate. Any zeros? \ans{No}
  \item So $k$ has \emph{no rational zeros}. But complete this table of values:
    \begin{center}
    \renewcommand{\arraystretch}{1.4}
    \begin{tabular}{|c|C{1.3cm}|C{1.3cm}|C{1.3cm}|C{1.3cm}|C{1.3cm}|}
    \hline
    $x$ & $-2$ & $-1$ & $0$ & $1$ & $2$ \\ \hline
    $k(x)$ & \ans{$-3$} & \ans{$1$} & \ans{$-1$} & \ans{$-3$} & \ans{$1$} \\ \hline
    \end{tabular}
    \end{center}
  \item The outputs change sign three times, so the graph must cross the $x$-axis
        \ans{$3$} times. Between which pairs of consecutive $x$-values?
        \ans{$(-2,-1)$, $(-1,0)$, and $(1,2)$}
  \item So $k$ has three \emph{real} zeros and \emph{no} rational ones. What kind of numbers must
        all three be? \ans{irrational} \; What does that tell you about what the Rational Root
        Theorem is --- and is not --- for? \ansline{It is a search tool for \emph{rational} zeros only. When a polynomial has none, the theorem returns nothing --- and that is not a contradiction, because irrational zeros were never in its reach.}
\end{enumerate}

\vspace{5pt}
\textbf{Part 4 --- Tomorrow's problem, today.} Solve $x^3 - x^2 + 4x - 4 = 0$.
\begin{enumerate}[label=(\alph*), leftmargin=1.8em, itemsep=6pt]
  \item Candidates: \ans{$\pm 1,\ \pm 2,\ \pm 4$} \; A zero that works: $r = $ \ans{$1$}
  \item Depressed polynomial: \ans{$x^2 + 4$} \; Try to solve it: $x^2 = $ \ans{$-4$}, so
        $x = $ \ans{$\pm\sqrt{-4}$ --- not a real number}
  \item How many \emph{real} solutions does the original equation have? \ans{$1$} \;
        How many solutions does a cubic ``want'' to have? \ans{$3$}
  \item Where did the other two go? Write your group's best guess --- Lesson 3.5 will tell you
        whether you were right. \ansline{They are imaginary: $\sqrt{-4} = 2i$, so the missing solutions are $x = \pm 2i$ --- a complex-conjugate pair, exactly as in Unit 2.}
\end{enumerate}
\end{tcolorbox}

\begin{teachernote}
\textbf{Tier placement.} Assign from the exit ticket of Lesson~3.3: the ``sign flip on $r$'' and
``arithmetic'' piles go to Tier~R, the ``stopped before factoring'' pile to Tier~A, and everyone else
to Tier~E. Tier~R problems 1--2 are monic on purpose, so the candidate list is just the divisors of
the constant term and students can practice the \emph{workflow} before the fraction bookkeeping
arrives in problem 3.

\vspace{3pt}
\textbf{Tier A problem 2 is the one to debrief with the whole class.} It is the lesson's pedagogical
core: the graph names all three zeros, the candidate list confirms them, and the fractional zero
$-\tfrac12$ is \emph{visible} as the crossing between $-1$ and $0$ before anybody divides anything.
Students who work top-to-bottom through twelve candidates get the right answer and miss the point;
ask them out loud how many divisions the graph saved. Problem 4(b) is the conceptual check --- a
candidate list with no winners is not a mistake, it is information.

\vspace{3pt}
\textbf{Tier E Part 1} is a genuine proof and worth the time; the argument is three lines and makes
the theorem stop feeling arbitrary. Part 2 inverts it, which is the skill Lesson~3.5 opens with
(building a polynomial from its zeros, A2.EI.6a). Part 3 is the honest limit of the tool, and Part 4
is the setup for the Fundamental Theorem of Algebra --- leave (d) as a conjecture rather than
telling them; groups that recall complex numbers from Unit~2 will get there themselves, and the ones
that do not will have a question waiting for tomorrow.

\vspace{3pt}
\textbf{If a group finishes early}, hand them the unit anchor $x^3-3x+2$ and ask them to confirm
that the Rational Root Theorem would have found both of its zeros with no hints at all --- the
candidate list is $\pm1,\pm2$, and $x=1$ and $x=-2$ are both on it. That closes the loop the class
has been walking since Lesson~3.0.
\end{teachernote}

\end{document}
